% Arquitectura (qué): topología del sistema, lista de nodos, dependencias, tipos de datos (mensajes), entradas y salidas de cada nodo, diagrama. NO explicar algoritmos, ni filtros, ni ecuaciones.
\section{Arquitectura general}
\label{sec:arquitectura-general}

El sistema se compone de un conjunto de nodos ROS~2 que se ejecutan sobre una Raspberry Pi~5,
orquestados por el archivo de lanzamiento \code{nav.launch.py} del paquete \texttt{nav\_bringup}.
Todos los nodos de procesamiento quedan contenidos en la misma unidad de cómputo y se comunican
exclusivamente a través del middleware DDS (\code{rmw\_cyclonedds}). La frontera de la
implementación física se ubica en la transmisión UART: el sistema genera y serializa los comandos
de intensidad, pero no incluye el receptor, la electrónica de potencia ni los actuadores.
Distintos nodos de visualización del paquete (\texttt{output\_viewer}) pueden ejecutarse
opcionalmente en una PC separada, suscribiéndose a los mismos tópicos a través de la red DDS. La
topología completa del sistema en su versión final integra percepción instantánea, conciencia
espacial y mapeo local, tal como se muestra en la \Cref{fig:arquitectura-v2}.

El Anexo~\ref{anexo:trazabilidad} muestra la correspondencia entre la arquitectura, las decisiones
metodológicas y el código fuente mediante un mapeo nodo → tópicos → archivos fuente.

\section{Comparación de la topología entre versiones}
\label{sec:arq-versiones}

Los diagramas de esta sección ofrecen una vista de conjunto de las tres configuraciones
incrementales definidas en la Sección~\ref{sec:intro-objetivo-tecnico}. Su propósito es mostrar
qué bloques se mantienen, se reemplazan o se agregan en cada versión; los tópicos, tipos de
mensaje e interfaces de cada nodo se detallan posteriormente en la
Sección~\ref{sec:arq-flujos-versiones}. La implementación conserva los modos \code{raw} y
\code{filtered}, seleccionados mediante \code{pipeline\_mode}, para ejecutar y comparar las
configuraciones alcanzadas en v0 y v1.

\subsection{Topología de v0: percepción directa}
La versión inicial utiliza exclusivamente la percepción instantánea: imagen de profundidad de la
cámara, conectada al pipeline mediante remapeos definidos en \texttt{nav\_bringup}, reducción a
\texttt{DepthGrid} en \texttt{depth\_grid\_encoder}, fusión trivial (passthrough), visualización de
la salida y serialización para su transmisión UART. La topología resultante se muestra en la
Figura~\ref{fig:arquitectura-v0}.

\begin{figure}[H]
  \centering
  \includegraphics[width=\textwidth]{arquitectura/Arquitectura-v0}
  \caption{Arquitectura v0.}
  \label{fig:arquitectura-v0}
\end{figure}

\subsection{Topología de v1: filtrado geométrico}
Sobre v0 se incorpora \texttt{depth\_obstacle\_filter} para suprimir suelo y techo, alineando la
nube a gravedad y conservando en la salida únicamente obstáculos relevantes para la marcha. Esta
etapa reemplaza la fuente directa de \texttt{DepthGrid}, pero conserva sin cambios los nodos de
generación háptica y transmisión. La topología resultante se muestra en la
Figura~\ref{fig:arquitectura-v1}.

\begin{figure}[H]
  \centering
  \includegraphics[width=\textwidth]{arquitectura/Arquitectura-v1}
  \caption{Arquitectura v1.}
  \label{fig:arquitectura-v1}
\end{figure}

\subsection{Topología de v2: memoria y riesgo contextual}
La tercera versión agrega memoria espacial mediante \texttt{local\_mapper} y una capa de riesgo
contextual con \texttt{spatial\_awareness}. Esta capa complementa la percepción frontal instantánea
con alertas de ocupación lateral o trasera, aprovechando información acumulada en el mapa local y
el estado de movimiento del usuario. La topología resultante se muestra en la
Figura~\ref{fig:arquitectura-v2}.

\begin{figure}[!htbp]
  \centering
  \includegraphics[width=\textwidth]{arquitectura/Arquitectura-v2}
  \caption{Arquitectura v2.}
  \label{fig:arquitectura-v2}
\end{figure}

\section{Infraestructura y adquisición de datos}

El sistema utiliza paquetes de soporte para definir los mensajes, los parámetros y los modos de
ejecución compartidos. Por separado, los tópicos publicados por el driver RealSense se remapean a
nombres internos estables. Así, los nodos de procesamiento no dependen de los nombres concretos
del driver y la configuración de la cámara puede cambiar sin modificar el resto del sistema.

\subsection{Paquetes de soporte e infraestructura}
\label{sec:arq-soporte}

El paquete \texttt{custom\_interfaces} concentra los tipos de mensajes específicos del sistema.
En la cadena de percepción frontal, \texttt{DepthGrid} modela una grilla compacta de celdas y
\texttt{DepthCellStats} encapsula la estadística de cada celda (media, mínimo, máximo y conteo
de puntos válidos). En la cadena de actuación, \texttt{HapticGrid} representa la salida final de
intensidades hápticas con dos arreglos explícitos: \texttt{intensities} y \texttt{active}. Para
la conciencia espacial, \texttt{DirectionalRisk} describe el riesgo por dirección
(activación, intensidad, distancia, velocidad de cierre y tiempo hasta colisión) y
\texttt{SpatialAwareness} agrupa los tres canales \texttt{left}, \texttt{right} y \texttt{rear}
junto con una bandera de validez global. Esta separación reduce el acoplamiento entre nodos y evita incorporar información de riesgo dentro de un mensaje destinado a percepción.

El paquete \texttt{nav\_bringup} centraliza la configuración de lanzamiento del sistema en
\code{nav.launch.py}. Este archivo expone parámetros para habilitar o deshabilitar componentes y
modos de ejecución, entre ellos: selección de pipeline de percepción
(\texttt{raw}/\texttt{filtered}/\texttt{none}), activación de \texttt{local\_mapper} y
\texttt{spatial\_awareness}, publicación en hardware y
herramientas de visualización y depuración. Estas opciones determinan qué nodos se inician y con
qué parámetros, mientras que la lógica de procesamiento permanece en los paquetes funcionales.

\subsection{Adquisición y normalización de tópicos}
\label{sec:arq-sensor}

El driver de la cámara, lanzado mediante \code{rs\_launch.py} del paquete \code
{realsense2\_camera}, publica los flujos del sensor bajo el espacio de nombres
\code{/camera/camera/}. En lo que sigue, para evitar repetición, ese prefijo se abrevia como
\code{RS/}. En esta interfaz conviene distinguir entre nombre de tópico y tipo de mensaje: por
ejemplo, \code{RS/depth/image\_rect\_raw}, \code{RS/color/image\_raw} y
\code{RS/aligned\_depth\_to\_color/image\_raw} son tópicos de imagen, todos transportados con
tipo \code{sensor\_msgs/Image}, mientras que \code{RS/depth/camera\_info} y
\code{RS/color/camera\_info} son tópicos distintos que publican parámetros intrínsecos con tipo
\code{sensor\_msgs/CameraInfo}. De forma análoga, las mediciones inerciales llegan por tópicos
específicos del driver y utilizan \code{sensor\_msgs/Imu} como tipo de mensaje.

El launcher \code{nav.launch.py} define remapeos explícitos para cada consumidor, de forma que la
lógica de procesamiento no dependa de los nombres publicados por el driver y se pueda identificar
qué flujo alimenta a cada nodo. Las correspondencias de referencia se indican en la
\Cref{tab:remapper}.

\begin{table}[htbp]
\centering
\begin{tabular}{ll}
\hline
\textbf{Tópico de origen (driver)} & \textbf{Tópico normalizado} \\
\hline
\texttt{RS/depth/image\_rect\_raw}  & \texttt{/sensors/depth/image} \\
\texttt{RS/depth/camera\_info}      & \texttt{/sensors/depth/camera\_info} \\
\texttt{RS/color/image\_raw}        & \texttt{/sensors/color/image} \\
\texttt{RS/color/camera\_info}      & \texttt{/sensors/color/camera\_info} \\
\texttt{RS/accel/sample}            & \texttt{/sensors/imu/accel} \\
\texttt{RS/gyro/sample}             & \texttt{/sensors/imu/gyro} \\
\hline
\end{tabular}
\caption{Correspondencias de tópicos implementadas mediante remapeos en \code{nav.launch.py}. En
esta tabla, \code{RS/} abrevia el prefijo \code{/camera/camera/}.}
\label{tab:remapper}
\end{table}

Además de estos remapeos, el código actual consume otros flujos del driver de forma directa.
\texttt{depth\_obstacle\_filter} utiliza \code{RS/depth/image\_rect\_raw},
\code{RS/depth/camera\_info} y \code{RS/accel/sample}; el filtro de Madgwick
(paquete \texttt{imu\_filter\_madgwick}) consume \code{RS/imu} para producir
\code{/imu/data\_filtered}; y \texttt{rtabmap\_odom} utiliza \code{RS/color/image\_raw},
\code{RS/color/camera\_info}, \code{RS/aligned\_depth\_to\_color/image\_raw} e
\code{/imu/data\_filtered}. La \Cref{tab:remapper} resume únicamente los remapeos
explícitos usados como interfaz de normalización, no la totalidad de tópicos de cámara presentes
en ejecución.

\section{Nodos y flujo de mensajes por versión}
\label{sec:arq-flujos-versiones}

Esta sección describe las interfaces que materializan las configuraciones comparadas en la
Sección~\ref{sec:arq-versiones}: qué nodo produce cada representación, qué mensajes intercambia y
qué componentes consumen la salida. Las tres versiones conservan un contrato de salida común,
pero modifican progresivamente la fuente de la información espacial. V0 obtiene la grilla
directamente de la profundidad; v1 la construye a partir de obstáculos filtrados y alineados con
gravedad; v2 mantiene la ruta frontal de v1 y agrega memoria local y riesgo contextual.

\subsection{Flujo de v0: profundidad y salida háptica}
\label{sec:arq-percepcion}

El nodo \code{depth\_to\_matrix} (paquete \texttt{depth\_grid\_encoder}) suscribe
\code{depth/image}, remapeado desde \code{RS/depth/image\_rect\_raw} en
\code{nav.launch.py}, y publica una representación compacta en
\code{/perception/depth\_grid} con tipo \code{custom\_interfaces/msg/DepthGrid}. El nodo
\code{haptic\_grid\_generator} consume esa interfaz y publica
\code{/perception/haptic\_grid} como \code{custom\_interfaces/msg/HapticGrid}; finalmente,
\code{grid\_to\_motor} recibe la matriz de intensidades para transmitirla por UART. Esta versión
establece la cadena funcional de referencia que conservan las versiones posteriores.

La Figura~\ref{fig:mensajes-v0} explicita el contrato de datos de esta ruta. Cada flecha indica el
nombre del tópico, el tipo de mensaje y su contenido relevante. El diagrama distingue así la
representación geométrica \texttt{DepthGrid}, la orden lógica \texttt{HapticGrid} y la interfaz
de transmisión. La construcción de la grilla y la conversión a intensidad se desarrollan en las
Secciones~\ref{sec:metodo-v0-depthgrid} y~\ref{sec:metodo-v0-intensidad}.

\begin{figure}[H]
  \centering
  \includegraphics[width=\textwidth]{mensajes/mensajes-v0}
  \caption{Mensajes funcionales de v0. La imagen de profundidad se reduce a
  \texttt{DepthGrid}, se transforma en una matriz lógica \texttt{HapticGrid} y se entrega a la
  interfaz de transmisión.}
  \label{fig:mensajes-v0}
\end{figure}

\subsection{Flujo de v1: filtrado y codificación de obstáculos}
\label{sec:arq-filtro-obstaculos}

El nodo \code{depth\_obstacle\_filter\_node}, del paquete
\texttt{depth\_obstacle\_filter}, suscribe la imagen de profundidad, sus parámetros
intrínsecos y el acelerómetro del driver. El nodo publica
\texttt{/depth\_obstacle\_filter/obstacle\_cloud}
como salida principal, y cuando se habilita la interfaz hacia mapeo local también
publica \texttt{/depth\_obstacle\_filter/free\_endpoints}. De forma adicional
emite \texttt{/depth\_obstacle\_filter/camera\_height} (\texttt{std\_msgs/Float32})
para informar la altura estimada de cámara respecto del suelo. En TF publica la
transformación dinámica
\texttt{gravity\_aligned\_frame}~$\to$~\texttt{camera\_depth\_optical\_frame}.
La nube de obstáculos es consumida por \code{obstacle\_grid\_encoder}, del paquete
\texttt{depth\_grid\_encoder}, que la codifica angularmente y publica
\code{/perception/depth\_grid}. De este modo, v0 y v1 ofrecen el mismo contrato
\texttt{DepthGrid} a las etapas posteriores, aunque lo construyen a partir de fuentes diferentes.

La Figura~\ref{fig:mensajes-v1} muestra ese reemplazo sin presentar v1 como una cadena
independiente. Las salidas funcionales hacia el codificador son la nube de obstáculos, expresada
en el marco alineado con gravedad, y la altura de cámara. El resultado vuelve a ser el mismo tipo
\texttt{DepthGrid} de v0, por lo que la generación háptica y la interfaz de salida no cambian.

\begin{figure}[H]
  \centering
  \includegraphics[width=\textwidth,trim=0 125 0 0,clip]{mensajes/mensajes-v1}
  \caption{Mensajes funcionales de v1. La imagen, los intrínsecos y el acelerómetro permiten
  obtener una nube de obstáculos alineada y codificarla conservando el contrato
  \texttt{DepthGrid}--\texttt{HapticGrid} de v0.}
  \label{fig:mensajes-v1}
\end{figure}

La separación entre el filtrado geométrico de v1 y la rama de odometría de v2 evita que la
percepción frontal dependa de la disponibilidad del mapa local. Los criterios de proyección,
alineación, clasificación y codificación se describen en las
Secciones~\ref{sec:metodo-v1-proyeccion} a~\ref{sec:metodo-v1-angular}.

\subsection{Flujo de v2: odometría, mapeo y fusión contextual}
\label{sec:arq-mapeo}

La versión v2 mantiene en la rama inferior de la Figura~\ref{fig:mensajes-v2} la percepción
frontal de v1 y agrega dos recorridos auxiliares. El primero obtiene \texttt{nav\_odom} mediante
odometría RGB-D: RTAB-Map consume color, sus intrínsecos, imagen de profundidad registrada en el
marco óptico de la cámara RGB y la orientación publicada por \texttt{imu\_filter\_madgwick}. El
segundo combina esa odometría con las nubes de evidencia ocupada y libre para construir
\texttt{OccupancyGrid}. A partir del mapa y del estado de movimiento, \texttt{spatial\_awareness}
publica el riesgo contextual que vuelve a \texttt{haptic\_grid\_generator}. Por lo tanto, la rama
superior no reemplaza la percepción instantánea: aporta una segunda entrada que se fusiona antes de
producir la única \texttt{HapticGrid} de salida.

\begin{figure}[H]
  \centering
  \includegraphics[width=\textwidth]{mensajes/mensajes-v2}
  \caption{Mensajes funcionales de v2. La ruta frontal se complementa con odometría, memoria
  local y riesgo contextual, que se fusiona antes de publicar la matriz háptica final. Se
  muestran únicamente los mensajes que intervienen en el procesamiento funcional principal.}
  \label{fig:mensajes-v2}
\end{figure}

El nodo \code{occupancy\_mapper\_node} (paquete \texttt{local\_mapper}) recibe las evidencias
ocupada y libre publicadas por \texttt{depth\_obstacle\_filter}, junto con
\texttt{nav\_odom} y la altura de cámara. Es responsable de mantener la memoria espacial y publica
\texttt{nav\_msgs/OccupancyGrid} en \texttt{/local\_mapper/occupancy\_grid}
con frame \texttt{odom}. También publica la TF dinámica
\texttt{odom}~$\to$~\texttt{gravity\_aligned\_frame} para mantener coherencia
de marcos entre los nodos de v2. El método de asociación temporal y actualización del mapa se
desarrolla en las Secciones~\ref{sec:metodo-v2-evidencia} y
\ref{sec:metodo-v2-ocupacion}.

El nodo \code{spatial\_awareness\_node} (paquete \texttt{spatial\_awareness}) suscribe la grilla
local, la odometría y el estado de apertura angular. Su responsabilidad es producir advertencias
contextuales para los sectores izquierdo, derecho y posterior. El resultado se publica en
\code{/spatial\_awareness/collision\_risk} con tipo
\code{custom\_interfaces/msg/SpatialAwareness}, que contiene un canal
\code{DirectionalRisk} por sector y un indicador \code{valid} para señalar
vigencia de sus entradas. El nodo
\code{haptic\_grid\_generator\_node} recibe este mensaje junto con
\code{/perception/depth\_grid} y publica una única salida consolidada en
\code{/perception/haptic\_grid} con tipo \code{custom\_interfaces/msg/HapticGrid}.
Los criterios de evaluación del riesgo y de fusión se desarrollan en las
Secciones~\ref{sec:metodo-v2-riesgo} y~\ref{sec:metodo-v2-fusion}.

\section{Interfaz lógica de salida y transmisión UART}
\label{sec:arq-hardware}

El nodo \code{grid\_to\_motor} (paquete \texttt{hardware\_manager}) suscribe
\code{/perception/haptic\_grid}; por lo tanto, recibe comandos de intensidad ya
calculados y no procesa distancias. Su responsabilidad es adaptar la matriz lógica al rango de
salida configurado y serializarla para su transmisión UART al microcontrolador. La implementación
del escalado, el formato de trama y el control de frecuencia se describe en el capítulo de
Implementación. La validación del receptor, de la electrónica de potencia y de los motores
vibrotáctiles queda fuera del alcance de esta tesis.

\section{Visualización y depuración}
\label{sec:arq-viz}

El nodo \code{depth\_grid\_heatmap} (paquete \texttt{output\_viewer})
puede ejecutarse opcionalmente en una PC conectada a la misma red DDS. Genera una visualización de
\code{/perception/haptic\_grid} y puede superponer el estado de
\code{/spatial\_awareness/collision\_risk} para inspección integrada de
percepción frontal y riesgo contextual. De este modo permite verificar la
estructura, orientación, actualización y coherencia de los comandos de intensidad
sin depender de la interfaz háptica física; no sustituye su validación perceptual.
Los nodos de percepción publican además interfaces opcionales de diagnóstico para inspeccionar en
RViz2 las representaciones intermedias. Estas salidas no intervienen en el flujo funcional.
El detalle exhaustivo de tópicos ROS~2, incluyendo productores, consumidores,
tipos de mensaje y correspondencia con el código fuente, se presenta en el
Anexo~\ref{anexo:trazabilidad}.
